\section{System Requirements}

MANAR uses the following operational and safety requirements to drive software design. Values in this section are \emph{design targets or simulated configuration limits}, not measured airframe capabilities.

\subsection{Operational Envelope and Performance Targets}
The system targets an operational flight envelope tailored for rugged wilderness and desert environments:
\begin{itemize}
    \item \textbf{Search Endurance Target ($SYS\text{-}G\text{-}02$):} At least 1.0 hour of continuous active search under a nominal payload and representative sensor duty cycle; physical endurance has not been measured.
    \item \textbf{Configured Altitude Limit ($SYS\text{-}G\text{-}01$):} The software test configuration clamps commanded altitude at $2000.0\text{ m}$. This is not evidence of climb, hover, regulatory, or propulsion capability at that altitude.
    \item \textbf{Speed Presets and Limits:} Four discrete velocity presets ($v_Q, v_A, v_I, v_H$) bounded by software velocity clamp $v_{\max}$ and altitude ceiling $H_{\max}$. Takeoff climb target is specified by $H_{\mathrm{launch}}$.
\end{itemize}

\subsection{Terrain-Adaptive Search Strategy}
Search path generation adapts to local terrain geometry:
\begin{itemize}
    \item \textbf{Open / Flat Terrain (implemented source):} Parallel-track rectangular sweeps parameterized by sector width $W$, height $H$, and row spacing $\Delta y$.
    \item \textbf{Mountainous / Steep Terrain (planned):} Terrain-aware contour following with a verified digital elevation model and sensor standoff constraint.
\end{itemize}

\subsection{Supervised-Autonomy Responsibility Allocation}
Table~\ref{tab:responsibility_allocation} delineates the functional boundaries between the human operator, deterministic control logic, and machine learning pipelines.

\begin{table}[t]
\caption{Supervised-Autonomy Responsibility Allocation}
\label{tab:responsibility_allocation}
\centering
\footnotesize
\begin{tabularx}{\columnwidth}{@{}lX@{}}
\toprule
\textbf{Layer} & \textbf{Core Responsibilities} \\
\midrule
\textbf{Human Operator} & Mission definition, geofencing, takeoff approval, real-time RTL/waypoint override, alert review, final probable-rescue classification, and guidance authorization. \\
\midrule
\textbf{Deterministic Core} & Implemented prototype navigation, route ordering, configuration clamps, plan locking, and RTL state changes; future autopilot integration remains planned. \\
\midrule
\textbf{Machine Learning} & Implemented raw visual candidate detections; planned acoustic/radar classification and temporal fusion. Machine-learning output is advisory and does not directly command actuators. \\
\bottomrule
\end{tabularx}
\end{table}

\subsection{Deterministic Contingency and Failsafe Rules}
The target flight system defines deterministic contingency responses. Only software mechanisms explicitly identified elsewhere as implemented should be read as present in the current artifact:
\begin{itemize}
    \item \textbf{Command-Link Loss ($SYS\text{-}R\text{-}01$, planned integration):} Hold, timeout, and RTL behavior requires autopilot and communication-link simulation before verification.
    \item \textbf{Energy Failsafe ($SYS\text{-}R\text{-}02$, simulated source):} Threshold-driven warning, RTL, and landing states can be exercised with the discrete battery test variable; return-energy sufficiency is not physically modeled.
    \item \textbf{GNSS Degradation (planned):} Attitude hold and local clearance sensing depend on future autopilot and sensor integration.
    \item \textbf{Environmental Limits (planned):} Wind, icing, and propulsion abort logic requires measured sensors and flight testing.
\end{itemize}
